% GENERATED FILE. Edit canonical lessons, then run npm run generate:books.
% canonical-source-hash: fnv1a64:9c30a41331237420
% canonical-lessons: KA-C398-pratibhatisu, KA-C398-horata, KA-C398-manavi, KA-C398-sanghatisu, KA-C398-karyakarta

\chapter{Protest and Struggle}
\label{ch:ka-protest-and-struggle}
\hlchaptermodality{398}

\begin{chapteropening}
\textbf{By the end of this chapter:} \emph{I can say to protest, a struggle, a petition, to mobilise and an activist.}

Complete the last lesson of chapter 398: I can say to protest, a struggle, a petition, to mobilise and an activist.
\end{chapteropening}

\section[\texorpdfstring{\kn{ಪ್ರತಿಭಟಿಸು} (pratibhaṭisu)}{ಪ್ರತಿಭಟಿಸು (pratibhaṭisu)}]{\kn{ಪ್ರತಿಭಟಿಸು} (pratibhaṭisu) — to protest}

\label{lesson:KA-C398-pratibhatisu}



\begin{quote}
Before the new one: say the Kannada for a community, then the Kannada for to unite.
\end{quote}

\subsection*{You'll want to know: \kn{ಪ್ರತಿಭಟಿಸು}}
\textbf{\kn{ಪ್ರತಿಭಟಿಸು}} — \emph{pratibhaṭisu} — \textquotedblleft{}to protest\textquotedblright{}.

\begin{quote}
\kn{ವಿದ್ಯಾರ್ಥಿಗಳು} \kn{ಶುಲ್ಕ} \kn{ಏರಿಕೆಯನ್ನು} \kn{ಪ್ರತಿಭಟಿಸಿದರು}.
\end{quote}

\begin{quote}
\emph{vidyārthigaḷu śulka ērikeyannu pratibhaṭisidaru.}
\end{quote}

\begin{quote}
The students protested the rise in fees.
\end{quote}

Past (I …): \kn{ಪ್ರತಿಭಟಿಸಿದೆ}, \emph{pratibhaṭiside}.

\kn{ಪ್ರತಿಭಟನೆ}, a protest.

\begin{cousinweb}
\textbf{\kn{ಪ್ರತಿಭಟಿಸು}} (\emph{pratibhaṭisu}) is built on Sanskrit \emph{pratibhaṭa}, \textquotedblleft{}a match for, rivalling\textquotedblright{} (Monier-Williams), with the Kannada verb ending -\emph{isu}: to stand up against.
\end{cousinweb}

\begin{grammarlens}[title={What you've built}]
One more word for giving reasons and opinions.
\end{grammarlens}

\subsection*{Your turn}
\begin{itemize}
\raggedright
  \item \emph{Say it:} \emph{pratibhaṭisu}
  \item \emph{Say it:} \emph{pratibhaṭisu}, once more
  \item \emph{Say it:} \emph{vidyārthigaḷu śulka ērikeyannu pratibhaṭisidaru.} — The students protested the rise in fees.
  \item \emph{Recall:} say the Kannada for a community, then the Kannada for to unite, then say \emph{pratibhaṭisu} again
\end{itemize}

\begin{tcolorbox}[breakable,colback=teal!4,colframe=teal!35!black,title={Before you move on}]
What does \kn{ಪ್ರತಿಭಟಿಸು} mean? (\textquotedblleft{}To protest\textquotedblright{}.) Say it once more.
\end{tcolorbox}
\section[\texorpdfstring{\kn{ಹೋರಾಟ} (hōrāṭa)}{ಹೋರಾಟ (hōrāṭa)}]{\kn{ಹೋರಾಟ} (hōrāṭa) — a struggle; a fight}

\label{lesson:KA-C398-horata}



\begin{quote}
Before the new one: say the Kannada for to unite, then the Kannada for to protest.
\end{quote}

\subsection*{You'll want to know: \kn{ಹೋರಾಟ}}
\textbf{\kn{ಹೋರಾಟ}} — \emph{hōrāṭa} — \textquotedblleft{}a struggle; a fight\textquotedblright{}.

\begin{quote}
\kn{ನೀರಿಗಾಗಿ} \kn{ಹಳ್ಳಿಯವರ} \kn{ಹೋರಾಟ} \kn{ಹತ್ತು} \kn{ವರ್ಷ} \kn{ನಡೆಯಿತು}.
\end{quote}

\begin{quote}
\emph{nīrigāgi haḷḷiyavara hōrāṭa hattu varṣa naḍeyitu.}
\end{quote}

\begin{quote}
The villagers' struggle for water went on for ten years.
\end{quote}

From \kn{ಹೋರಾಡು}, to fight.

\begin{cousinweb}
\textbf{\kn{ಹೋರಾಟ}} (\emph{hōrāṭa}) is a native Dravidian compound, the noun beside \kn{ಹೋರಾಡು}: \emph{pōr}, \textquotedblleft{}to fight, to strive\textquotedblright{} (DEDR 4540), whose \emph{p-} became \emph{h-}, and \emph{āṭa}, \textquotedblleft{}moving, playing\textquotedblright{}, from \emph{āḍu} (DEDR 347).
\end{cousinweb}

\begin{grammarlens}[title={What you've built}]
One more word for giving reasons and opinions.
\end{grammarlens}

\subsection*{Your turn}
\begin{itemize}
\raggedright
  \item \emph{Say it:} \emph{hōrāṭa}
  \item \emph{Say it:} \emph{hōrāṭa}, once more
  \item \emph{Say it:} \emph{nīrigāgi haḷḷiyavara hōrāṭa hattu varṣa naḍeyitu.} — The villagers' struggle for water went on for ten years.
  \item \emph{Recall:} say the Kannada for to unite, then the Kannada for to protest, then say \emph{hōrāṭa} again
\end{itemize}

\begin{tcolorbox}[breakable,colback=teal!4,colframe=teal!35!black,title={Before you move on}]
What does \kn{ಹೋರಾಟ} mean? (\textquotedblleft{}A struggle; a fight\textquotedblright{}.) Say it once more.
\end{tcolorbox}
\section[\texorpdfstring{\kn{ಮನವಿ} (manavi)}{ಮನವಿ (manavi)}]{\kn{ಮನವಿ} (manavi) — a petition; an appeal}

\label{lesson:KA-C398-manavi}



\begin{quote}
Before the new one: say the Kannada for to protest, then the Kannada for a struggle.
\end{quote}

\subsection*{You'll want to know: \kn{ಮನವಿ}}
\textbf{\kn{ಮನವಿ}} — \emph{manavi} — \textquotedblleft{}a petition; an appeal\textquotedblright{}.

\begin{quote}
\kn{ಜನರು} \kn{ಮಂತ್ರಿಗೆ} \kn{ಒಂದು} \kn{ಮನವಿ} \kn{ಕೊಟ್ಟರು}.
\end{quote}

\begin{quote}
\emph{janaru mantrige ondu manavi koṭṭaru.}
\end{quote}

\begin{quote}
The people gave the minister a petition.
\end{quote}

\kn{ಮನವಿ} \kn{ಮಾಡು}, to appeal.

\begin{cousinweb}
\textbf{\kn{ಮನವಿ}} (\emph{manavi}) is a Kannada word whose origin is not settled: neither DEDR nor CDIAL lists it.
\end{cousinweb}

\begin{grammarlens}[title={What you've built}]
One more word for giving reasons and opinions.
\end{grammarlens}

\subsection*{Your turn}
\begin{itemize}
\raggedright
  \item \emph{Say it:} \emph{manavi}
  \item \emph{Say it:} \emph{manavi}, once more
  \item \emph{Say it:} \emph{janaru mantrige ondu manavi koṭṭaru.} — The people gave the minister a petition.
  \item \emph{Recall:} say the Kannada for to protest, then the Kannada for a struggle, then say \emph{manavi} again
\end{itemize}

\begin{tcolorbox}[breakable,colback=teal!4,colframe=teal!35!black,title={Before you move on}]
What does \kn{ಮನವಿ} mean? (\textquotedblleft{}A petition; an appeal\textquotedblright{}.) Say it once more.
\end{tcolorbox}
\section[\texorpdfstring{\kn{ಸಂಘಟಿಸು} (saṅghaṭisu)}{ಸಂಘಟಿಸು (saṅghaṭisu)}]{\kn{ಸಂಘಟಿಸು} (saṅghaṭisu) — to mobilise; to organise}

\label{lesson:KA-C398-sanghatisu}



\begin{quote}
Before the new one: say the Kannada for a struggle, then the Kannada for a petition.
\end{quote}

\subsection*{You'll want to know: \kn{ಸಂಘಟಿಸು}}
\textbf{\kn{ಸಂಘಟಿಸು}} — \emph{saṅghaṭisu} — \textquotedblleft{}to mobilise; to organise\textquotedblright{}.

\begin{quote}
\kn{ಅವಳು} \kn{ಕೆಲಸಗಾರರನ್ನು} \kn{ಸಂಘಟಿಸಿದಳು}.
\end{quote}

\begin{quote}
\emph{avaḷu kelasagārarannu saṅghaṭisidaḷu.}
\end{quote}

\begin{quote}
She organised the workers.
\end{quote}

Past (I …): \kn{ಸಂಘಟಿಸಿದೆ}, \emph{saṅghaṭiside}.

\kn{ಸಂಘಟನೆ}, an organisation.

\begin{cousinweb}
\textbf{\kn{ಸಂಘಟಿಸು}} (\emph{saṅghaṭisu}) is built on Sanskrit \emph{saṃghaṭana}, \textquotedblleft{}union, junction\textquotedblright{} (Monier-Williams), with the Kannada verb ending -\emph{isu}.
\end{cousinweb}

\begin{grammarlens}[title={What you've built}]
One more word for giving reasons and opinions.
\end{grammarlens}

\subsection*{Your turn}
\begin{itemize}
\raggedright
  \item \emph{Say it:} \emph{saṅghaṭisu}
  \item \emph{Say it:} \emph{saṅghaṭisu}, once more
  \item \emph{Say it:} \emph{avaḷu kelasagārarannu saṅghaṭisidaḷu.} — She organised the workers.
  \item \emph{Recall:} say the Kannada for a struggle, then the Kannada for a petition, then say \emph{saṅghaṭisu} again
\end{itemize}

\begin{tcolorbox}[breakable,colback=teal!4,colframe=teal!35!black,title={Before you move on}]
What does \kn{ಸಂಘಟಿಸು} mean? (\textquotedblleft{}To mobilise; to organise\textquotedblright{}.) Say it once more.
\end{tcolorbox}
\section[\texorpdfstring{\kn{ಕಾರ್ಯಕರ್ತ} (kāryakarta)}{ಕಾರ್ಯಕರ್ತ (kāryakarta)}]{\kn{ಕಾರ್ಯಕರ್ತ} (kāryakarta) — an activist}

\label{lesson:KA-C398-karyakarta}



\begin{quote}
Before the new one: say the Kannada for a petition, then the Kannada for to mobilise.
\end{quote}

\subsection*{You'll want to know: \kn{ಕಾರ್ಯಕರ್ತ}}
\textbf{\kn{ಕಾರ್ಯಕರ್ತ}} — \emph{kāryakarta} — \textquotedblleft{}an activist\textquotedblright{}.

\begin{quote}
\kn{ಕಾರ್ಯಕರ್ತರು} \kn{ಬೀದಿಯಲ್ಲಿ} \kn{ಮರ} \kn{ನೆಟ್ಟರು}.
\end{quote}

\begin{quote}
\emph{kāryakartaru bīdiyalli mara neṭṭaru.}
\end{quote}

\begin{quote}
The activists planted trees in the street.
\end{quote}

Plural: \kn{ಕಾರ್ಯಕರ್ತರು}, \emph{kāryakartaru}.

Built from the Sanskrit words for work and for a doer.

\begin{cousinweb}
\textbf{\kn{ಕಾರ್ಯಕರ್ತ}} (\emph{kāryakarta}) is borrowed from Sanskrit \emph{kāryakartṛ}, \textquotedblleft{}one who works in the interest of another\textquotedblright{} (Monier-Williams).
\end{cousinweb}

\begin{grammarlens}[title={What you've built}]
One more word for giving reasons and opinions.
\end{grammarlens}

\subsection*{Your turn}
\begin{itemize}
\raggedright
  \item \emph{Say it:} \emph{kāryakarta}
  \item \emph{Say it:} \emph{kāryakarta}, once more
  \item \emph{Say it:} \emph{kāryakartaru bīdiyalli mara neṭṭaru.} — The activists planted trees in the street.
  \item \emph{Recall:} say the Kannada for a petition, then the Kannada for to mobilise, then say \emph{kāryakarta} again
\end{itemize}

\begin{tcolorbox}[breakable,colback=teal!4,colframe=teal!35!black,title={Before you move on}]
What does \kn{ಕಾರ್ಯಕರ್ತ} mean? (\textquotedblleft{}An activist\textquotedblright{}.) Say it once more.
\end{tcolorbox}
